\begin{textsample}{POS Dim 6 – ms – Score 10.00 – ms/ms\_em\_35.txt}  \label{ex:f6_pos_054}
4.2.
Eixos \textbf{Estruturantes} A Portaria MEC n. 1.432/2018 corrobora o parágrafo 2º do art. 12 das DCNEM e estabelece a organização dos Itinerários \textbf{Formativos} a partir de quatro eixos \textbf{estruturantes} : Investigação Científica, Processos Criativos, \textbf{Mediação} e Intervenção Sociocultural e \textbf{Empreendedorismo}.
Esses eixos são complementares e auxiliarão os estudantes a desenvolverem habilidades necessárias para uma formação integral, que promova situações de aprendizagem realmente significativas para o seu crescimento intelectual e emocional.
Os estudantes terão a possibilidade de percorrer um ou, preferencialmente, todos os eixos \textbf{estruturantes} dos Itinerários \textbf{Formativos}, quais sejam : Investigação Científica — os estudantes podem participar da realização de uma pesquisa científica, compreendida como procedimento privilegiado e integrador de áreas e componentes curriculares.
O processo pressupõe a identificação de uma dúvida, questão ou \textbf{problema} ; o levantamento, formulação e teste de hipóteses ; a seleção de informações e de fontes confiáveis ; a interpretação, elaboração e uso ético das informações coletadas ; a identificação de como utilizar os conhecimentos gerados para solucionar \textbf{problemas} diversos ; e a comunicação de conclusões com a utilização de diferentes linguagens.
Processos Criativos — os estudantes podem realizar projetos criativos por meio da utilização e integração de diferentes linguagens, manifestações sensoriais, vivências artísticas, culturais, midiáticas e científicas aplicadas.
O processo pressupõe a identificação e o aprofundamento de um tema ou \textbf{problema}, que orientará a posterior elaboração, apresentação e difusão de uma ação, produto, protótipo, modelo ou solução criativa, tais como obras e espetáculos artísticos e culturais, campanhas e peças de comunicação, programas, \textbf{aplicativos}, jogos, robôs, circuitos, dentre outros produtos analógicos e \textbf{digitais}. \textbf{Mediação} e Intervenção Sociocultural — privilegia -se o envolvimento dos estudantes em campos de atuação da vida pública, por meio do seu engajamento em projetos de mobilização e intervenção sociocultural e \textbf{ambiental} que os levem a promover transformações positivas para a comunidade.
O processo pressupõe o diagnóstico da realidade sobre a qual se pretende atuar, incluindo a busca de dados oficiais e a escuta da comunidade local ; a ampliação de conhecimentos sobre o \textbf{problema} a ser enfrentado ; o planejamento, execução e avaliação de uma ação social e/ou \textbf{ambiental} que responda às necessidades e interesses do contexto ; a superação de situações de estranheza, resistência, conflitos interculturais e possíveis obstáculos com necessários ajustes de rota. \textbf{Empreendedorismo} — os estudantes são estimulados a criar empreendimentos pessoais ou produtivos articulados com seus projetos de vida que fortaleçam a sua atuação como protagonistas da sua própria trajetória.
Para isso, busca desenvolver autonomia, foco e determinação para que consigam planejar e conquistar objetivos pessoais ou criar empreendimentos voltados à geração de renda via oferta de produtos e serviços, com ou sem uso de \textbf{tecnologias}.
O processo pressupõe a identificação de potenciais, desafios, interesses e aspirações pessoais ; a análise do contexto externo, inclusive em relação ao mundo do trabalho ; a elaboração de um projeto pessoal ou produtivo ; a realização de ações-piloto para testagem e aprimoramento do projeto elaborado ; o desenvolvimento ou aprimoramento do projeto de vida dos estudantes.
Portanto, ao escolher os Itinerários \textbf{Formativos} que desejam percorrer ao longo do Ensino Médio, os estudantes poderão ter acesso a diferentes possibilidades de aprendizagens, com vistas a fortalecer o protagonismo, a criatividade, a atuação social, a proposição e a resolução de conflitos.

% matched lemmas: ambiental, aplicativo, digital, empreendedorismo, estruturante, formativo, mediação, problema, tecnologia
\end{textsample}
